\section{System and Threat Model}
\label{sec:threat}

\subsection{Entities}
\emph{Issuer $I$.} The tenant's authentication service, or an issuer operated by the CDN, holds the signing key $sk$. $sk$ is an ML-DSA-44/65, SLH-DSA-128s/f or hybrid Ed25519\,+\,ML-DSA-44 key, identified by a key id $\mathit{kid}$. \emph{Root service.} An origin object path from which signed epoch roots and period bundles are served. It is itself cacheable through the CDN's shield tier. \emph{Edges} $E_1,\dots,E_E$ hold $pk$, a root cache and a revocation set. \emph{Clients} authenticate to $I$ (for example at a paywall), obtain a token and then issue a manifest request and $K$ segment requests, normally sticky to one edge.

\subsection{Why Not Simply HMAC?}
HMAC-SHA256 tokens are secure against quantum adversaries. Their weakness is that every verifier can also mint. (i) A single edge compromise (a memory disclosure, an insider, a misconfiguration, partner hardware) yields a key that mints tokens for every edge and every path the key covers. (ii) A content owner using several CDNs must share the minting secret with each vendor, and each vendor can mint tokens that the others accept. (iii) There is no issuer accountability. Public-key tokens give edges only $pk$. We do not overstate the benefit. A compromised edge can already serve anything it caches without checking tokens, so asymmetric tokens remove only \emph{lateral} minting: tokens accepted by other edges, other CDNs or other tenants. Table~\ref{tab:trust} in Section~\ref{sec:results} quantifies this blast radius for every baseline.

\subsection{Adversary}
The adversary $\mathcal{A}$ (i) controls the network between clients and edges and can replay observed tokens; (ii) may compromise \emph{any subset} of edges and shields, learning their full state and injecting roots into their caches; (iii) is a legitimate client of any tenant and obtains arbitrarily many tokens for claims of its choice, while leaf salts remain issuer-chosen; and (iv) is quantum. It can run Shor's algorithm against any public key and generic quantum search against the hash function. $\mathcal{A}$ does not obtain the issuer's $sk$. Issuer compromise is handled by key rotation via $\mathit{kid}$. Theft and replay of bearer tokens are unchanged from the status quo and are mitigated by the same means: IP or ASN claims, short validity, and optionally a DPoP-style key-hash claim~\cite{rfc9449}.

\subsection{Goals and Requirements}
\textbf{G1 (unforgeability):} no efficient (quantum) $\mathcal{A}$ makes an honest edge accept claims that $I$ never issued or has revoked, even if all other edges are compromised. \textbf{G2 (edge non-minting):} edges hold no secret. \textbf{G3 (bounded validity):} an edge rejects a token whose expiry exceeds $T_{\text{start}}(e)+\Delta_e+\vmax$ (plus a clock-skew tolerance $\tau$). \textbf{G4 (sibling privacy):} an authentication path reveals nothing about other clients' claims. \textbf{G5 (availability):} an edge that cannot obtain a root fails closed.

Operational requirements are: \textbf{R1}, tokens fit URL and cookie budgets with margin; \textbf{R2}, per-request edge cost close to HMAC; \textbf{R3}, no minting secret at edges; \textbf{R4}, bounded issuance latency; \textbf{R5}, revocation; \textbf{R6}, negligible dissemination bandwidth and rare root-miss waits; and \textbf{R7}, a clear security reduction.
